% Part: methods
% Chapter: induction
% Section: strong-induction

\documentclass[../../../include/open-logic-section]{subfiles}

\begin{document}

\olfileid{mth}{ind}{str}

\olsection{প্রবল আরোহ}

আগে আলোচিত আরোহের নীতিতে আমরা $P(0)$ প্রমাণ করি,
আরও প্রমাণ করি যে $P(k)$ হলে $P(k+1)$। দ্বিতীয় অংশে
$P(k)$ সত্য ধরে সেই অনুমান থেকে $P(k+1)$ প্রমাণ করি।
সমতুল্যভাবে, $P(k-1)$ ধরে তা থেকে $P(k)$-ও প্রমাণ
করতে পারি। আসল কথা হলো, যেকোনো সংখ্যা থেকে তার
উত্তরসূরিতে যাওয়ার যুক্তিটি যেন প্রয়োগ করতে পারি:
যেকোনো সংখ্যার পূর্বসূরির জন্য দাবিটি সত্য ধরে সংখ্যাটির
জন্য দাবি প্রমাণ করতে পারি।
% BN-SRC-839 TARGET-CORRECTION-BEGIN
\par\noindent\textbf{পাঠ-সংশোধন।} এই অধ্যায়ে সংখ্যার পরিমাণায়ন স্বাভাবিক সংখ্যায়।
পূর্বসূরি-রূপে $P(k-1)$ থেকে $P(k)$ ধাপটি $k\geq1$-এর
জন্য; শূন্যের কোনও স্বাভাবিক পূর্বসূরি নেই। প্রবল আরোহের
পূর্বধারণায় সব $l<k$ বলতে স্বাভাবিক $l$-গুলি বোঝায়।
তাই $k=0$-তে ওই সার্বিক পূর্বধারণা শূন্যক্ষেত্রে সত্য,
আর সব $k$-এর শর্তাধীন ধাপ নিজেই $P(0)$ দেয়। শুধুমাত্র
ধনাত্মক $k$-এর জন্য ধাপ প্রমাণ করলে আলাদা শূন্য-ভিত্তি
এখনও দরকার।
% BN-SRC-839 TARGET-CORRECTION-END

আরোহের নীতির একটি রূপভেদে $k$-এর পূর্বসূরি $k-1$-এর
জন্যই শুধু দাবিটি সত্য ধরা হয় না; $k$-এর চেয়ে ছোট
সব সংখ্যার জন্য সত্য ধরা হয়। এই অনুমান থেকে $k$-এর
জন্য দাবি প্রতিষ্ঠা করি। এতেও সব~$n \in \Nat$-এর জন্য
$P(n)$ পাওয়া যায়। কারণ $P(0)$ প্রতিষ্ঠা করলেই $1$-এর
চেয়ে ছোট সব সংখ্যার জন্য~$P$ সত্য প্রতিষ্ঠিত হয়। আর
যদি জানি যে সব $l<k$-এর জন্য $P(l)$ হলে $P(k)$, তবে
বিশেষভাবে $k=1$-এর জন্যও তা জানি। তাই $P(1)$ সিদ্ধান্ত
করি। এখন $P(0)$ ও $P(1)$—অর্থাৎ সব $l<2$-এর জন্য
$P(l)$—প্রমাণিত। পাশাপাশি ``সব $l<2$-এর জন্য $P(l)$
হলে $P(2)$'' শর্তাধীন বক্তব্যটিও আমাদের আছে। তাই~$P(2)$
পাই। এভাবেই পরের সংখ্যাগুলির জন্যও এগোই।

আসলে ``সব $k$-এর জন্য, সব $l<k$-এর জন্য $P(l)$ হলে
$P(k)$'' সাধারণ শর্তাধীন বক্তব্যটি প্রতিষ্ঠা করতে পারলে
$P(0)$ আর আলাদা করে প্রতিষ্ঠা করতে হয় না; ওই বক্তব্য
থেকেই তা অনুসৃত। মনে রাখো, ``সব $l<k$-এর জন্য $P(l)$''
ধরনের সাধারণ বক্তব্য সত্য যখন কোনো $l<k$-ই নেই। এটি
শূন্যক্ষেত্রে সত্য পরিমাণায়নের উদাহরণ: কোনো $A$ না থাকলে
``সব $A$-ই $B$'' সত্য; কোনো $x$ যদি~$!A(x)$ পরিতৃপ্ত না
করে, তবে $\lforall[x][(!A(x) \lif !B(x))]$ সত্য।
এখানে আনুষ্ঠানিক রূপটি হবে ``$\lforall[l][(l < k \lif P(l))]$'';
কোনো $l<k$ না থাকলে এটিও সত্য। $k=0$ হলে ঠিক এই
অবস্থাই ঘটে: কোনো $l<0$ নেই। তাই
% BN-SRC-484: the vacuous antecedent has P(l), not P(0).
``সব $l<0$-এর জন্য $P(l)$'' সত্য, $P$ যা-ই হোক।
সুতরাং ``সব $l<k$-এর জন্য $P(l)$ হলে $P(k)$'' প্রমাণ
করলেই স্বয়ংক্রিয়ভাবে~$P(0)$ প্রতিষ্ঠিত হয়।

যখন $k$-এর জন্য দাবি প্রমাণে শুধু $k-1$-এর দাবির
উপর নির্ভর করা যায় না, কিন্তু এক বা একাধিক $l<k$-এর
জন্য সত্যতার অনুমান দরকার হতে পারে, তখন এই রূপভেদ
কাজে লাগে।

\end{document}
